\subsubsection{Contextual Transformation Networks (CTN)}

\gls{ctn} \cite{ctn_2021} is an online, bi-level \gls{cl} method comprising two
components: a backbone network and a controller network. The backbone extracts
feature representations from the input, while the controller applies a learned
linear transformation to those features. This decoupled design enables
task-specific feature modulation whilst sharing a common backbone, in contrast
to methods such as \cite{ll_lwf_2017, ll_ewc_2017} that apply identical feature
representations across all tasks. \gls{ctn} operates under the \gls{til}
assumption, requiring task identity to be available at both training and
inference time in order to select the appropriate per-task transformation
parameters.

The controller network is implemented as a \gls{mlp} that maps a task embedding
vector to per-task \gls{film} parameters \cite{2018_FiLM},
$\{\gamma_t, \beta_t\} \in \mathbb{R}^d$, where $d$ denotes the dimensionality
of the backbone feature vector $\hat{h}(x) \in \mathbb{R}^d$. These parameters
modulate $\hat{h}(x)$ element-wise as follows:
%
\begin{equation}\label{eqn:ctn_FiLM}
    \tilde{h}(x;t) = \frac{\gamma_t}{\|\gamma_t\|_2} \odot \hat{h}(x)
    + \frac{\beta_t}{\|\beta_t\|_2}
\end{equation}
%
where $\odot$ denotes element-wise multiplication. The $\ell_2$ normalisation
of $\gamma_t$ and $\beta_t$ constrains the scale of the transformation,
stabilising training by preventing the controller outputs from growing
unboundedly. The final task-conditioned feature vector is obtained via a
residual connection:
%
\begin{equation}\label{eqn:ctn_final_features}
    h(x;t) = \tilde{h}(x;t) + \hat{h}(x)
\end{equation}

The backbone parameters $\phi$ and controller parameters $\theta$ are optimised
in a bi-level scheme: $\phi$ is updated using episodic memory-augmented training
data, whilst $\theta$ is updated on a separate semantic memory stream. These two
update steps are interleaved at each training iteration.

\paragraph{\gls{cl} Approach}

\gls{ctn} mitigates \gls{cf} through two complementary mechanisms: experience
replay and feature-level regularisation.

For experience replay, \gls{ctn} maintains two memory buffers: an episodic
memory (\texttt{memx}, \texttt{memy}) and a semantic memory (\texttt{valx},
\texttt{valy}). Each incoming minibatch is partitioned such that one sample is
diverted to the semantic memory, whilst the remaining samples are used for
backbone training and episodic replay. Backbone and controller updates are
therefore driven by distinct but synchronised sample streams: $\phi$ is
optimised using episodic memory-augmented training data, whilst $\theta$ is
updated using samples drawn from the semantic memory. Because the controller
trains on samples withheld from backbone updates, it encourages the backbone to
produce features that generalise beyond the immediate training distribution.
Both memories are bounded in capacity and partitioned evenly across tasks; once
full, they operate as ring buffers in which incoming samples overwrite the
oldest entries. This bounded capacity is consistent with the online \gls{cl}
setting, in which the full training data stream is not retained.

For output-distribution regularisation, \gls{ctn} stores replay exemplars from each
task and, at task transitions, caches softened class-probability targets for those
stored samples (rather than intermediate feature vectors). During subsequent tasks,
training uses replay classification plus a distillation penalty that matches current
and cached prior-task output distributions, in the spirit of \cite{ll_lwf_2017}. The
inner-loop training objective is:
%
\begin{equation}
\mathcal{L}^{tr}
= \mathcal{L}_{\mathrm{CE}}\!\left(\hat{\mathbf{y}}, \mathbf{y}\right)
+ \mathcal{L}_{\mathrm{replay}}\!\left(\hat{\mathbf{y}}_{\mathrm{rep}}, \mathbf{y}_{\mathrm{rep}}\right)
+ \lambda \,
D_{\mathrm{KL}}\!\left(
\log \pi\!\left(\frac{\hat{\mathbf{y}}_{\mathrm{rep}}}{\tau}\right)
\,\Big\|\,
\pi\left(\dfrac{\hat{\mathbf{y}}^{\mathrm{old}}_{\mathrm{rep}}}{\tau}\right)
\right),
\end{equation}
%
where $\pi(\cdot)$ is softmax, $\hat{\mathbf{y}}_{\mathrm{rep}}$ are current logits on
replayed samples, $\pi(\hat{\mathbf{y}}^{\,\mathrm{old}}_{\mathrm{rep}}/\tau)$ are cached softened targets
computed from the previous-task model at transition time, $\tau>0$ is the temperature,
and $\lambda$ controls distillation strength.
Thus, \gls{kl} regularisation is applied in output space (task-local predictive
distributions), indirectly constraining representation drift and mitigating forgetting.

\paragraph{Task Embeddings}

The controller must learn a distinct feature transformation for each task,
requiring a task representation as input. A natural choice is a one-hot encoding
of the task index; however, this is sparse and encodes no information about
inter-task relationships. Instead, \gls{ctn} learns a dense task-embedding
vector, updated via backpropagation alongside $\{\phi, \theta\}$, which serves
as input to the controller when generating the \gls{film} parameters
$\{\gamma_t, \beta_t\}$. This allows the task representation to become
increasingly informative as training progresses, and in principle enables the
controller to exploit task similarity when modulating backbone features.
